Let \(y=h+\mathbb Z_0\) and write \(T=T_{b_0}(s_0)\).  Since \(\mathbb B(s_0)=\{b_0\}\), the definitions in \(\text{def:boundary-limit-sets}\) reduce the zeroth-order union to
\[
 \mathcal K_h^{(0)}
 =\{\theta_{b_0}(s_0):T\cap\mathcal A_\gamma(y)\ne\varnothing\}.
 \tag{1}
\]
Thus it is the displayed singleton when the intersection is nonempty and is empty otherwise.  Under the same singleton-branch condition, the first-order definition gives, with \(J=D\theta_{b_0}(s_0)\),
\[
 \mathcal K_h^{(1)}
 =\{Ju:u\in T\cap\mathcal A_\gamma(y)\}
 =J\bigl(T\cap\mathcal A_\gamma(y)\bigr)
 =\mathcal P_\gamma(y;b_0,s_0),
 \tag{2}
\]
where the final equality is exactly \(\text{def:regular-product-score-image}\).  Equations (1)--(2) are also the singleton-branch specialization of \(\text{thm:boundary-random-set-limit}\).

If \(s_0\) lies in the ambient interior of \(\mathcal S_{b_0}\), then some open ball about \(s_0\) is contained in the cell.  Every \(u\in\mathbb R^d\) is therefore a feasible first-order direction, so
\[
 T_{b_0}(s_0)=\mathbb R^d.
 \tag{3}
\]
Writing \(u=(u_1,u_2)\), \(J=(J_1,J_2)\), and using the product definition \(\mathcal A_\gamma(y)=\mathcal A_{1,\gamma}(y_1)\times\mathcal A_{2,\gamma}(y_2)\), linearity gives both inclusions in
\[
 \begin{aligned}
 J\mathcal A_\gamma(y)
 &=\{J_1u_1+J_2u_2:
       u_e\in\mathcal A_{e,\gamma}(y_e),\ e=1,2\}\\
 &=J_1\mathcal A_{1,\gamma}(y_1)
   \oplus J_2\mathcal A_{2,\gamma}(y_2).
 \end{aligned}
 \tag{4}
\]
If \(\mathcal A_{e,\gamma}(y_e)=y_e+E_{e,\gamma}\), another application of linearity yields
\[
 J\mathcal A_\gamma(y)
 =J_1y_1+J_2y_2+J_1E_{1,\gamma}\oplus J_2E_{2,\gamma}
 =Jy+J_1E_{1,\gamma}\oplus J_2E_{2,\gamma}.
 \tag{5}
\]
This is a Minkowski sum of two linear ellipsoid images; no identity makes it a joint Wald ellipsoid in general.

Finally, every branch target has the typed form
\[
 \theta_b(s)=(\alpha_1(s),\alpha_2(s),
                    \alpha_1(s)-\alpha_2(s)).
 \tag{6}
\]
Differentiating (6), which is permitted by \(\text{ass:branch-smoothness}\), shows that the third row of \(J\) is the first row minus the second.  Hence \(\operatorname{rank}(J)\le2\).  No lower-rank or full-rank premise was used in (1)--(5), proving the proposition without a rank assumption.